\section{Matrix extraction and Fourier decoding}
\label{sec:matrix}

We convert high acceptance of the matrix test into agreement between two
strategies for the sparsely projected questions $U,O$.  Each strategy
uses only its own question and a bounded number of rows of matrix advice.
The resulting agreement probability is positive independently of the tuple
length $k$ and the parity instance.  The table functions enter only through
their folding and sharing identities and the acceptance bound, so the proof
applies to any family with these properties; the restored functions of every vertex
labeling are examples.  We state the precise assumptions below, after
fixing the parameters and defining the advice experiment.

\subsection{The inverse shortcode theorem}

The structural input describes functions that often keep their value
under a uniform rank-one perturbation.  Such a function agrees with an
affine evaluation on a set obtained by fixing a bounded number of linear
combinations of rows and columns.

Write $\operatorname{Mat}_{\ell,m}=\Hom(\F_2^m,\F_2^\ell)$.
A \emph{row and column slice} is a nonempty set specified by equations
\begin{equation}
 \mathcal S=
 \{M:d_i^\top M=s_i^\top\ (1\le i\le n_r),\quad
          Mq_j=t_j\ (1\le j\le n_c)\},
 \label{eq:shortcode-slice}
\end{equation}
where $d_i\in\F_2^\ell$, $s_i\in\F_2^m$, $q_j\in\F_2^m$, and
$t_j\in\F_2^\ell$.  The counts $n_r,n_c$ need not be the ranks of their
respective systems.

\begin{theorem}[Equality soundness for degree-two shortcode]
\label{thm:shortcode}
For every $\eta\in(0,1)$ there are $\alpha\in(0,1]$, an integer
$r_{\rm s}\ge1$, and an integer $\ell_0$, depending only on $\eta$,
with the following property.  For every $\ell\ge\ell_0$ there is
$m_0=m_0(\eta,\ell)$ such that, for every $m\ge m_0$ and every function
$f:\operatorname{Mat}_{\ell,m}\to\F_2^\ell$, if
\begin{equation}
 \Pr_{M,a,l}\bigl[f(M)=f(M+a l^\top)\bigr]\ge\eta,
 \label{eq:shortcode-acceptance}
\end{equation}
then some slice $\mathcal S$ as in Equation~\eqref{eq:shortcode-slice},
with $n_r,n_c\le r_{\rm s}$, and some $z\in\F_2^m$, $u\in\F_2^\ell$
satisfy
\begin{equation}
 \Pr_{M\mid\mathcal S}[f(M)=Mz+u]\ge\alpha.
 \label{eq:shortcode-conclusion}
\end{equation}
Here $M$, $a\in\F_2^\ell$, and $l\in\F_2^m$ are independent and
uniform, and $M\mid\mathcal S$ is uniform on the slice.  No folding
condition is imposed on $f$.
\end{theorem}

\begin{proof}
We derive this consequence from the Grassmann expansion theorem of
Khot--Minzer--Safra \cite[Theorem~1.12 in the 2018 revision]{KMS23},
using the matrix chart of Barak--Kothari--Steurer
\cite[Definition~3.2 and Lemmas~3.3--3.7]{BKS18}.

The Grassmann graph on an ambient binary space $E$ has the
$\ell$-dimensional subspaces as vertices; two vertices are adjacent when
their intersection has dimension $\ell-1$.  The \emph{retention} of a
nonempty vertex set is the probability that a uniform vertex in the set
has a uniform neighbor still in the set.  The cited theorem says that a
set of retention at least a fixed $\zeta>0$ has positive relative
density in some interval
\[
 \{L:A_0\subseteq L\subseteq B_0\},
 \qquad \dim A_0+\operatorname{codim}_E B_0\le r.
\]
The density bound and $r$ depend only on $\zeta$.  The theorem holds for
all sufficiently large $\ell$, and then all sufficiently large
$\dim E$.

Equality acceptance is the average of the retention probabilities of the
nonempty fibers $f^{-1}(y)$ under the matrix step, weighted by their
uniform masses.  Some fiber $S$ therefore has retention at least $\eta$.
The probability that either perturbation factor is zero is
\[
 p_0=2^{-\ell}+2^{-m}-2^{-\ell-m}.
\]
Choose the dimension thresholds so that $p_0\le\eta/2$.  Conditional on
both factors being nonzero, the retention of $S$ is at least
$(\eta-p_0)/(1-p_0)\ge\eta/2$.

For each matrix put
\[
 L_M=\{(x,M^\top x):x\in\F_2^\ell\}
 \subseteq E:=\F_2^\ell\oplus\F_2^m.
\]
The map $M\mapsto L_M$ is a bijection onto the chart $\mathcal C$ of
subspaces whose projection onto the first summand is invertible.  Since
\[
 \dim(L_M\cap L_N)=\ell-\rank(M-N),
\]
the neighbors within the chart correspond exactly to nonzero rank-one
matrix differences.  Each chart vertex has
$(2^\ell-1)(2^m-1)$ such neighbors, whereas its full Grassmann degree is
$(2^\ell-1)(2^{m+1}-2)$: choose a hyperplane in the vertex and then a
one-dimensional extension in the quotient by that hyperplane, other than
the vertex itself.  Thus exactly
half of its neighbors lie in the chart.  Over $\F_2$, every nonzero
rank-one matrix has a unique factorization $a l^\top$ with nonzero
factors.  The lifted set $\widetilde S=\{L_M:M\in S\}$ consequently has
Grassmann retention at least $\eta/4$.

Apply the expansion theorem with $\zeta=\eta/4$, obtaining density
$\alpha$ on an interval with
$\dim A_0+\operatorname{codim}_E B_0\le r_{\rm s}$.  Intersecting
this interval with $\mathcal C$ only increases the relative density of
$\widetilde S$, since $\widetilde S\subseteq\mathcal C$, and the
intersection is nonempty.  A basis $(d_i,s_i)$ of $A_0$ specifies
$A_0\subseteq L_M$ by the equations $d_i^\top M=s_i^\top$.
For the standard dot product on $E$, a basis $(t_j,q_j)$ of
$B_0^\perp$ specifies $L_M\subseteq B_0$ by
$Mq_j=t_j$.  Their simultaneous solution set is a slice of the required
form, with at most $r_{\rm s}$ equations of either kind.  On at least an
$\alpha$ fraction of this slice, $f$ equals the value defining $S$.
This proves the conclusion with $z=0$ and $u$ equal to that value.

The constants $\alpha,r_{\rm s}$ were chosen using only $\eta$.
Since $\dim E=\ell+m$, after fixing $\ell$ the remaining dimension
threshold is absorbed into $m_0(\eta,\ell)$.  No folding condition was
used.  The fiber argument is the equality consequence in
\cite[Lemma~2.3]{BKS18}; numbered references to that paper use its
arXiv version~1.
\end{proof}

The constants and dimension thresholds in Theorem~\ref{thm:shortcode}
can be chosen effectively, with rational $\alpha$ when $\eta$ is rational.
This follows from the quantitative bound in
\cite[Theorem~2.13 and Section~2.4]{KMS23} and the comparison estimates
in \cite[Lemmas~2.7--2.8]{KMS23}, with theorem numbers again referring to
the 2018 revision.

We now fix the parameters needed below.  Choose $p_*>0$ and the associated
$r_*$ from Lemma~\ref{lem:latent}.  Set
\begin{equation}
 \eta=2^{-r_*}/16,
 \qquad
 2^{-(\ell-r_{\rm s})}<\alpha/8,
 \qquad
 2^{-\ell}<\eta/4,
 \label{eq:matrix-parameters}
\end{equation}
where $\alpha,r_{\rm s}$ first come from Theorem~\ref{thm:shortcode},
and $\ell$ is then chosen sufficiently large for that theorem, for the
latent lemma, and for the displayed inequalities.  Obtain the latent data
$(\mathcal V,K,C,\mu)$ at this value of $\ell$, and fix a linear splitting
$\mathcal V=K\oplus W$.  All these objects are fixed before $k$ grows.
In particular, $\ell>r_{\rm s}$.  We write a table as $P=(M,T)$, with
$M:L_U\to K$ and $T:L_U\to W$.

\subsection{The advice experiment and the decoding statement}

We first specify the local inputs for two strategies whose answer agreement
will contradict the outer-game gap.  This experiment is defined for every positive integer $k$
and every $0<\beta<1$.  Let $\mathcal Q_k$ consist of all equation
questions $U\in\mathcal E^k$ and all permitted projected questions $O$
obtained from them by the coordinate projections of Section~\ref{sec:game}.
The questions $U,O$ and the projection
$\bar\pi:L_U\to L_O$ have the law in Section~\ref{sec:game}: independently
in each block, an equation is retained with probability $1-\beta$, and a
uniformly chosen variable position is retained with probability $\beta$.

\begin{definition}[Projected advice experiment]
\label{def:projected-advice}
Sample $U,O,\pi$, then an independent uniform map
$A:K\to\F_2^{r_{\rm s}}$ and independent uniform maps
\[
 T_O:L_O\to W,\qquad M_O:L_O\to K.
\]
Put $S_O=AM_O$ and send the following inputs to the two provers:
\begin{equation}
 \begin{aligned}
  &\text{$U$-prover:}&& (U,A,T_O\bar\pi,S_O\bar\pi),\\
  &\text{$O$-prover:}&& (O,A,T_O,S_O).
 \end{aligned}
 \label{eq:projected-advice}
\end{equation}
The map $M_O$ beyond the stated advice is hidden.  Each prover returns an
actual answer, that is, a vector of first bit one in its homogeneous
answer space.  They succeed when their answers $a_U,a_O$ satisfy
$\bar\pi a_U=a_O$.
\end{definition}

The $U$-message contains the maps
$(T_U,S_U)=(T_O\bar\pi,S_O\bar\pi)$, obtained by padding maps sampled
on $O$.  We will first find useful advice under a different law: an
unrestricted uniform table $(M,T)$ on $U$, with $S_U=AM$.  Later we
compare the full padded table $(M_O\bar\pi,T_O\bar\pi)$ with this
unrestricted table after averaging over the projection choices.  This
comparison is the bridge from the extraction argument to the experiment.

\begin{proposition}[Matrix decoding]
\label{prop:matrix-decoding}
Choose the latent and shortcode parameters in the order specified by
Equation~\eqref{eq:matrix-parameters}, fixing $\mathcal V=K\oplus W$
before $k$.  There are a constant $\gamma>0$ and an integer $k_0\ge8$, depending only on
these fixed parameters, with the following property.

For every cube $k\ge k_0$ and every parity instance, suppose a function
\[
 F_Q:\Hom(L_Q,\mathcal V)\longrightarrow K
\]
is given for every $Q\in\mathcal Q_k$.  Assume folding for every such
$Q$, every $P\in\Hom(L_Q,\mathcal V)$, and every $c\in K$, and assume
sharing for every permitted projection $\bar\pi:L_U\to L_O$ and every
$P'\in\Hom(L_O,\mathcal V)$:
\begin{equation}
 \begin{aligned}
 F_Q(P+c e_0^\top)&=F_Q(P)+c,\\
 F_U(P'\bar\pi)&=F_O(P').
 \end{aligned}
 \label{eq:matrix-family-identities}
\end{equation}
If this family passes the table test of Equation~\eqref{eq:matrix-test}
with probability at least $0.99$, then the projected advice experiment
with $\beta=k^{-2/3}$ has local strategies whose probability of full
answer agreement is at least $\gamma$.
\end{proposition}

Every labeling of the vertices of $\mathcal G$ supplies such a family:
its restored functions satisfy Equation~\eqref{eq:matrix-family-identities}
by Equations~\eqref{eq:table-folding} and~\eqref{eq:key-sharing}.  The
abstract formulation records that no other property of the vertex
construction enters the decoding argument.  The family is fixed when the
two response rules are defined and may be built into those rules.

The next two subsections work for any positive integer $k$, with a fixed
family satisfying Equation~\eqref{eq:matrix-family-identities}.  They use
unrestricted uniform tables on $U$ and do not yet restrict $\beta$.
We specialize $k$ and $\beta$ when comparing that law with the projected
experiment.

\subsection{Comparing the two matrix perturbations}

The uniform matrix perturbation has the rank-indexed Fourier spectrum
computed in \cite[Lemma~2.10 in the 2018 revision]{KMS23}. We compare
that spectrum with the noise supplied by the latent gadget.

\begin{lemma}[Spectral comparison]
\label{lem:matrix-spectral}
If the table family passes Equation~\eqref{eq:matrix-test} with probability at
least $0.99$, then
\begin{equation}
 \E_{U,T}\Pr_{M,a,l}
   [F_U(M,T)=F_U(M+a l^\top,T)]
 \ge 0.9\,2^{-r_*},
 \label{eq:surrogate-acceptance}
\end{equation}
where $a$ is uniform in $K$ and the other matrix and vector samples are
uniform in their respective spaces.  Consequently at least $10\eta$ of
the pairs $(U,T)$ have the inner probability at least $4\eta$.
\end{lemma}

\begin{proof}
For $\sigma\in K^*$ define the sign function
$f_{U,\sigma}(P)=(-1)^{\sigma(F_U(P))}$.  A frequency on the matrix
domain is a tuple $g=(g_i)_{i=0}^{m-1}$ of functionals in $\mathcal V^*$;
its character at $P$ is $(-1)^{\sum_i g_i(Pe_i)}$.  All Fourier
coefficients use uniform probability measure.  Translation by $a l^\top$
multiplies this character by $(-1)^{\sum_i l_i g_i(a)}$.  Averaging over
uniform $l$ therefore gives the eigenvalue
\[
 \Lambda_g=\Pr_{a\sim\mu}[g_i(a)=0\text{ for every }i].
\]
It lies in $[0,1]$.  By Lemma~\ref{lem:latent}, it is at most $7/8$
whenever
\[
 \dim\Span\{g_i|_K:0\le i<m\}\ge r_*.
\]
Indeed the span of the $g_i$ is a subspace of $\mathcal V^*$ to which
the lemma applies.

Expanding equality by a uniform output character, and then using Fourier
orthogonality, expresses the test acceptance as
\[
 \E_{U,\sigma}\sum_g
      \abs{\widehat f_{U,\sigma}(g)}^2\Lambda_g.
\]
Parseval gives total squared mass one for each sign function.  If $B$ is
the average squared mass on restricted ranks at least $r_*$, then
$0.99\le1-B/8$, so $B\le0.08$.  In particular, at least $0.9$ of the
average mass lies on ranks below $r_*$.

When $a$ is instead uniform in $K$, the eigenvalue at $g$ is
$2^{-\dim\Span\{g_i|_K\}}$.  The preceding mass thus gives the lower
bound in Equation~\eqref{eq:surrogate-acceptance}.  This perturbation
leaves $T$ fixed, so the expression there is precisely the average of the
ordinary shortcode equality tests at fixed $(U,T)$.  Its lower bound is
$14.4\eta$.  If $q$ is the probability that the inner acceptance is at
least $4\eta$, then the average is at most $4\eta+q$.  Thus
$q\ge10.4\eta\ge10\eta$.
\end{proof}

\subsection{Obtaining many useful row fibers}

The inverse theorem gives one structured slice on each qualifying
$(U,T)$ fiber.  The $U$-prover's information about the $K$-valued matrix
is the pair $(A,AM)$.  We need a positive probability that these observed rows determine
a useful slice.  We will also identify the short set of valid answers
supplied by that slice.  The erasure argument below proves the probability
bound, and folding then makes the answer set nonempty.

For this extraction step, sample $U$ as before and take independent uniform
maps $T:L_U\to W$, $A:K\to\F_2^{r_{\rm s}}$, and $M:L_U\to K$.
For fixed $U,T,A$ and a map
$S_0:L_U\to\F_2^{r_{\rm s}}$, the set $\{M:AM=S_0\}$ is its
\emph{row fiber}.  We call the row advice $(U,T,A,S_0)$ \emph{good} if there exist
at most $r_{\rm s}$ column specifications and a target $Mz+u$ such that
the simultaneous solution set of the column specifications and $AM=S_0$
is nonempty and $F_U(M,T)=Mz+u$ with probability at least $\alpha/2$
under the uniform measure on this set.

This is an analytic condition on the row fiber.  The target vector $z$
need not initially be an actual answer.  The fixed column values allow
$z$ to vary modulo the span of their column directions, with a corresponding
change in $u$.  Thus the slice naturally supplies a coset of possible
target vectors.  Its valid part will be the answer set sampled by the
$U$-prover.

\begin{lemma}[Good row advice and its short answer set]
\label{lem:good-row-advice}
For all sufficiently large $k$, a table family with matrix-test acceptance at
least $0.99$ satisfies
\begin{equation}
 \Pr_{U,T,A,M}[(U,T,A,AM)\text{ is good}]
 \ge g_0,
 \qquad
 g_0=10\eta^2\,2^{-\ell r_{\rm s}}>0.
 \label{eq:good-advice-mass}
\end{equation}
The lower bound and the required largeness of $k$ are independent of the
parity instance and the family.

For every positive integer $k$ and every family satisfying
Equation~\eqref{eq:matrix-family-identities}, one can select deterministically
for each good row advice at most $r_{\rm s}$ column specifications
$Mq_j=t_j$, with $q_j\in L_U$ and $t_j\in K$, and a target $Mz+u$,
with $z\in L_U$, $u\in K$, and $e_0^\top z=1$, satisfying the definition
of goodness.
Writing $Z=\Span\{q_j\}$, the chosen witness determines the answer set
\begin{equation}
 \mathcal A=(z+Z)\cap\{v\in L_U:e_0^\top v=1\},
 \qquad 1\le |\mathcal A|\le 2^{r_{\rm s}}.
 \label{eq:row-answer-set}
\end{equation}
The witness and $\mathcal A$ depend only on $(U,T,A,S_0)$ and the fixed
table family.  The slice and its target retain the agreement guarantee in
the definition of goodness.
\end{lemma}

\begin{proof}
Fix $(U,T)$ for which the surrogate acceptance is at least $4\eta$, and
write $f(M)=F_U(M,T)$.  Let $B$ be the union, in the matrix domain, of all
good row fibers, over every possible $A,S_0$.  Suppose that its uniform
measure is less than $\eta$.  Change $f$ at the entries of $B$, replacing
each by an independent uniform element of $K$.  Each matrix receives one
replacement, even when several good row fibers overlap there.  Both endpoints of the
shortcode step are uniform matrices.  Consequently every such modification
changes acceptance by at most $2\Pr[B]$, and the modified function has
acceptance at least $2\eta$.

We show that the random replacements can be chosen so that they contribute
less than $\alpha/4$ of the size of every simultaneous row and column slice
as matches with every target $Mz+u$.  Put $r=r_{\rm s}$.  A nonempty
slice with at most $r$ row and $r$ column specifications has homogeneous
space
\[
 \Hom\bigl(\F_2^m/\Span\{q_j\},\ \bigcap_i\ker d_i^\top\bigr).
\]
Its size $n_{\mathcal S}$ is therefore at least
$2^{(\ell-r)(m-r)}$.  Padding the lists of specifications with zero
equations, the number of slice--target descriptions is at most
$2^{(2r+1)(\ell+m)}$.

For a fixed description, count only matches at randomized entries and
assign a deterministic zero contribution to the other entries.  These are
independent variables in $[0,1]$, with total mean at most
$2^{-\ell}n_{\mathcal S}<\alpha n_{\mathcal S}/8$.  Hoeffding's
inequality \cite[Theorem~2, equation~(2.6)]{Hoeffding1963} bounds the
probability of at least
$\alpha n_{\mathcal S}/4$ such matches by
$\exp(-\alpha^2 n_{\mathcal S}/32)$.  The union bound over all
descriptions is thus at most
\[
 \exp\left((2r+1)(\ell+m)\log2
       -\frac{\alpha^2}{32}2^{(\ell-r)(m-r)}\right),
\]
which tends to zero as $m=1+2k$ grows.  Fix replacements with the asserted
property.

Consider any slice of the allowed type.  If its row fiber was good, every
entry of that fiber belongs to $B$, and its modified agreement is less
than $\alpha/4$.  If the row fiber was not good, the original agreement
with each target on the slice was less than $\alpha/2$.  The unmodified
matches form a subset of those original matches, and the randomized
matches add less than $\alpha/4$.  The modified agreement is then less
than $3\alpha/4$.  Padding shorter row systems by zeros shows that these
two cases cover every slice in Theorem~\ref{thm:shortcode}.  That theorem
now contradicts the modified function's acceptance at least $2\eta$.
Its applicability does not require the modification to preserve folding.

It follows that $\Pr[B]\ge\eta$.  For each $M\in B$, at least one
matrix $A$ has a good row fiber containing $M$.  The uniform choice of $A$
hits that particular matrix with probability $2^{-\ell r_{\rm s}}$.
Hence, at this fixed $(U,T)$, the probability of good data is at least
$\eta2^{-\ell r_{\rm s}}$.  Lemma~\ref{lem:matrix-spectral} supplies
at least $10\eta$ of such pairs $(U,T)$, proving
Equation~\eqref{eq:good-advice-mass}.  All counting and concentration
bounds depended only on the fixed parameters and $m$, as required.

The remaining assertion does not use high test acceptance or largeness of
$k$.  Fix any positive integer $k$, a family satisfying
Equation~\eqref{eq:matrix-family-identities}, and any good row advice.
Fix a witness, write its column specifications as $Mq_j=t_j$, and put
$Z=\Span\{q_j\}$.  Nonemptiness of the slice makes the prescribed values
consistent: they determine a linear map $t:Z\to K$ with $t(w)=Mw$
throughout the slice.  Replacing $z$ by $z+w$ and $u$ by $u+t(w)$ for
$w\in Z$ therefore preserves the target.  It suffices to show that the
coset $z+Z$ contains a vector of first bit one.

Otherwise $z$ and every vector in $Z$ have first bit zero.  Put $H=\ker A$.
For each $h\in H$, translation $M\mapsto M+h e_0^\top$ preserves the
row fiber, the column specifications, and the target.  By folding it
translates $F_U(M,T)$ by $h$.  The action is free, and in each of its
$H$-orbits at most one matrix can satisfy the target equality.  The
agreement is at most
\[
 |H|^{-1}=2^{-\ell+\rank A}
 \le2^{-(\ell-r_{\rm s})}<\alpha/8,
\]
contradicting goodness.  Choose $w\in Z$ for which $z+w$ has first bit one
and make the corresponding changes to $z$ and $u$.  They preserve the
original agreement fraction.  Choose a witness and such a representative
using a fixed order on the finite possibilities.  This choice uses only
the row advice and the fixed table family.  Its vector $z$ lies in
$\mathcal A$, so the set is nonempty; since $\dim Z\le r_{\rm s}$,
its size is at most $2^{r_{\rm s}}$.
\end{proof}

\subsection{Transferring the chosen slice}

The comparison with sparsely projected questions follows the smooth
outer-game method of Khot--Safra~\cite[Section~3 and Lemma~3.1]{KS13},
with shared linear advice as in Khot--Minzer--Safra~\cite[Sections~3.3--3.4]{KMS17}.
Lemma~\ref{lem:good-row-advice} gives a positive probability of row advice
that determines a slice and a short set of valid answers under unrestricted
sampling on $U$.  In the actual experiment, the maps are sampled on $O$
and padded to $U$.  We compare the complete padded table with the
unrestricted table, including the affine intercept, with projection choices
averaged out, and then transfer the chosen slice to $O$.

From now on take $k$ along integer cubes larger than one and set
$\beta=k^{-2/3}$ in Definition~\ref{def:projected-advice}.

The choice $\beta=k^{-2/3}$ serves two purposes: $k\beta^2\to0$ makes
the padded table law approach the unrestricted law in
Lemma~\ref{lem:projected-tv}, while $k\beta\to\infty$ makes the expected
number of singleton coordinates diverge.  Once the advice spaces are
fixed, each singleton has a uniformly positive probability that every
visible advice row has zero coefficient on its variable bit.
Lemma~\ref{lem:clean} uses these coordinates to bound agreement.

\begin{lemma}[The padded matrix is asymptotically uniform]
\label{lem:projected-tv}
Conditional on every $U$, the distribution of
\[
 (M,T)=(M_O\bar\pi,T_O\bar\pi)
\]
has total variation distance $o(1)$ from the unrestricted uniform
distribution on $\Hom(L_U,K)\times\Hom(L_U,W)$, as $k$ tends to
infinity along cubes.  The bound is uniform in $U$ and the parity instance.
The projection choices are averaged out in this comparison.
\end{lemma}

\begin{proof}
Combine the two components into a $\mathcal V$-valued table and put
$n=|\mathcal V|$.  On an unprojected equation block the two slopes are
independent uniform elements of $\mathcal V$.  On a singleton block the
padded pair of slopes is an equal mixture of
\[
 (a,0),\qquad (0,a),\qquad (a,a),
 \qquad a\text{ uniform in }\mathcal V.
\]
Call these pair laws $P_0$ and $Q_0$, respectively.  The likelihood ratio
$dQ_0/dP_0$ is $n$ at $(0,0)$, $n/3$ at the other $3(n-1)$ points on
the three displayed lines, and zero elsewhere.  Thus
\[
 \chi^2(Q_0\Vert P_0)
 :=\E_{P_0}\left[\left(\frac{dQ_0}{dP_0}-1\right)^2\right]
 =\frac{n-1}{3}.
\]
The one-block mixture $(1-\beta)P_0+\beta Q_0$ consequently has
chi-squared divergence $\beta^2(n-1)/3$ from $P_0$.  Blocks are sampled
independently, including when their question IDs coincide.  Multiplying
their likelihood ratios gives divergence
\begin{equation}
 \left(1+\frac{n-1}{3}\beta^2\right)^k-1=o(1),
 \label{eq:projected-chi-square}
\end{equation}
because $n$ is fixed and $k\beta^2=k^{-1/3}$ tends to zero.  Total
variation is at most one half the square root of this expression.

For a projection onto the third bit of equation $e$, the padded intercept
acquires the term $b_e a$, where $a$ is that singleton's slope.  The
original intercept is independent and uniform in $\mathcal V$.  Adding
the sum of these terms leaves it uniform and independent even after all
projection choices and slopes are given.  The preceding slope estimate
therefore proves the assertion for the entire table.  Neither the slope
laws nor the bound in Equation~\eqref{eq:projected-chi-square} depend on
the equation IDs or their right-hand sides.
\end{proof}

Fix good $U$-advice for a moment, with its chosen witness and answer set
$\mathcal A$ from Lemma~\ref{lem:good-row-advice}.  For a permitted
projection put $z'=\bar\pi z$ and $Z'=\bar\pi Z$.  Since $\bar\pi$
preserves the first coordinate,
\begin{equation}
 \bar\pi\mathcal A
 =(z'+Z')\cap\{v\in L_O:e_0^\top v=1\}.
 \label{eq:projected-answer-set}
\end{equation}
Indeed, every vector in the coset on the right has a preimage in $z+Z$,
and the preimage has the same first coordinate.  We will show that the
$O$-prover's private Fourier sample reaches this projected answer set.

The certifying slice does not automatically survive projection.  For
example, a column condition $Mq=t$ becomes $M_O\bar\pi q=t$; it is
inconsistent if $\bar\pi q=0$ but $t\ne0$.  More generally, projected
column directions may acquire new dependencies.  The next argument uses
the comparison of the full table laws to find a positive probability of samples
where the chosen slice remains nonempty and retains enough agreement.

For this analysis, write
$D=(U,O,\pi,A,T_O,S_O)$ for the joint data that determine both prover
messages and the projection.  These data do not specify $M_O$ beyond the
constraint $AM_O=S_O$.

\begin{lemma}[Transfer of a chosen witness]
\label{lem:projected-witness}
Suppose the matrix-test acceptance is at least $0.99$.  For all sufficiently
large cubes $k$, a set of probability at least
\begin{equation}
 \gamma_1=\frac{\alpha g_0}{8}\,2^{-\ell r_{\rm s}}>0
 \label{eq:projected-witness-mass}
\end{equation}
of the joint data $D$ in
Definition~\ref{def:projected-advice} has the following properties.
The padded $U$-advice is good, with the chosen answer set $\mathcal A$ and its
certifying witness.  Writing $z'=\bar\pi z$, that witness transfers to a
nonempty slice
\[
 AM_O=S_O,\qquad M_O\bar\pi q_j=t_j,
\]
on which
\begin{equation}
 \Pr\bigl[F_O(M_O,T_O)=M_Oz'+u\mid
           AM_O=S_O,\ M_O\bar\pi q_j=t_j\text{ for all }j\bigr]
 \ge\alpha/4.
 \label{eq:transferred-witness}
\end{equation}
The probability in this equation is taken over a uniform matrix in the
specified slice.
\end{lemma}

\begin{proof}
First use unrestricted sampling of $U,A,M,T$, setting $S_0=AM$.  Let
$\mathcal J$ be the event that $(U,T,A,S_0)$ is good and $M$ satisfies
its chosen column specifications.  Let $\mathcal M$ be equality with
the chosen target, declaring this event false off good data.  Within a
good row fiber the nonempty column slice
imposes at most $\ell r_{\rm s}$ additional scalar linear equations.
Hence Lemma~\ref{lem:good-row-advice} gives
\[
 \Pr[\mathcal J]\ge g_0\,2^{-\ell r_{\rm s}},
 \qquad
 \Pr[\mathcal J\cap\mathcal M]\ge(\alpha/2)\Pr[\mathcal J].
\]
It follows that
\begin{equation}
 \Pr[\mathcal J\cap\mathcal M]
       -(\alpha/4)\Pr[\mathcal J]
 \ge\frac{\alpha g_0}{4}\,2^{-\ell r_{\rm s}}.
 \label{eq:witness-positive-difference}
\end{equation}

Both events in this expression depend only on $(U,A,M,T)$.  Adding the
independent $A$ to Lemma~\ref{lem:projected-tv} preserves its variation
bound.  We may therefore compare the two event probabilities with their
values in the padded experiment.  If the total variation distance is
$v_k=o(1)$, the difference in Equation~\eqref{eq:witness-positive-difference}
changes by at most $(1+\alpha/4)v_k$.  For all sufficiently large $k$
the padded difference is at least $\gamma_1$.

Now condition on $D$.  The remaining $M_O$ is uniform on the row fiber
$AM_O=S_O$.  The padded
$U$-advice, and therefore its chosen witness and answer set, are fixed.
On good data let $p_D$ be the conditional probability of the projected
column slice, and, when $p_D>0$, let $q_D$ be its conditional equality
fraction.  Put $p_D=0$ on data that are not good, and put $q_D=0$ when
$p_D=0$.  Sharing gives
$F_U(M_O\bar\pi,T_O\bar\pi)=F_O(M_O,T_O)$, so the padded difference
is exactly
\[
 \E_D\bigl[p_D(q_D-\alpha/4)\bigr]\ge\gamma_1.
\]
The integrand is nonpositive outside the event consisting of goodness,
$p_D>0$, and $q_D\ge\alpha/4$, and it is at most one everywhere.
That event thus has probability at least $\gamma_1$ and has precisely the
asserted properties.  The law comparison applied to events depending only
on $(U,A,M,T)$; conditioning on $D$ here only identifies where their
positive difference occurs in the projected experiment.
\end{proof}

\subsection{Decoding the row advice}

Extracting an outer strategy from Fourier mass follows the approach of
H\aa stad~\cite[Section~1.1]{Hastad01} and the Fourier-list method of
Khot--Safra~\cite[Section~1.1]{KS13}. Here the projected prover samples
squared Fourier mass on its own row fiber.  The transferred slice will
show that this private sample reaches the projected answer set
$\bar\pi\mathcal A$ with positive probability.

We now define both strategies from their separate messages.  The joint
data and the transferred witness serve only in the analysis of their
agreement; the $O$-prover does not use the $U$-prover's witness or answer set.

The $U$-prover uses its visible row advice to make the deterministic choice
in Lemma~\ref{lem:good-row-advice}.  When the advice is good, it samples
$a_U$ uniformly from $\mathcal A$.
This is a nonempty set of at most $2^{r_{\rm s}}$ actual answers.  On
other inputs the prover returns any actual answer.

The $O$-prover starts from its own row fiber.  Given $(O,A,T_O,S_O)$, put
$H=\ker A$ and choose, by a fixed rule, a representative $M_0$ satisfying
$AM_0=S_O$.  Such a representative exists
on every input in the experiment.  Write the row fiber as $M_0+X$, where
$X=\Hom(L_O,H)$.  Sample uniform $\sigma\in K^*$ and form the sign
function
\[
 f_\sigma(N)=(-1)^{\sigma(F_O(M_0+N,T_O))},\qquad N\in X.
\]
Sample a Fourier frequency $\Phi\in X^*$ with probability
$|\widehat f_\sigma(\Phi)|^2$.  Identify
$X^*=H^*\otimes L_O$, with the pairing defined on pure tensors by
\[
 \langle\varphi\otimes v,N\rangle=\varphi(Nv).
\]
For $h\in H$, contraction sends $\Phi$ to a vector $\Phi(h)\in L_O$.
If $\sigma|_H\ne0$, choose by a fixed rule an $h\in H$ with $\sigma(h)=1$ and
return $\Phi(h)$.  If the restriction is trivial, return any actual
answer.

This rule always returns an actual answer at frequencies of positive
mass.  Indeed the folding identity in
Equation~\eqref{eq:matrix-family-identities} gives
\[
 f_\sigma(N+h e_0^\top)=-f_\sigma(N)
 \qquad\text{when }\sigma(h)=1.
\]
Changing variables in a Fourier coefficient shows that its being nonzero
requires
$(-1)^{\langle\Phi,h e_0^\top\rangle}=-1$, equivalently
$e_0^\top\Phi(h)=1$.

Figure~\ref{fig:advice} summarizes the information and private choices of
the two decoders.  The transferred slice will be used only to analyze
their agreement probability.
\begin{figure}[tbp]
\centering
\begin{tikzpicture}[
  x=1cm,y=1cm,
  every node/.style={font=\small},
  box/.style={draw=linkblue!80!black,rounded corners=2pt,
    text width=6.0cm,align=center,inner sep=8pt},
  input/.style={box,minimum height=16mm},
  rule/.style={box,minimum height=35mm},
  output/.style={box,minimum height=10mm},
  flow/.style={-{Latex[length=2.1mm]},thick,draw=black!70}
]
\node[box,text width=13.3cm,minimum height=12mm] (sample) at (3.7,0)
  {\textbf{Referee}\\[2pt]
   Sample $U,O,\bar\pi,A,T_O,M_O$; set $S_O=AM_O$.
   Send only the inputs below.};

\node[input] (uinput) at (0,-1.9)
  {\textbf{Equation decoder}\\
   $(U,A,T_U,S_U)$\\
   $T_U=T_O\bar\pi,\quad S_U=S_O\bar\pi$};
\node[input] (oinput) at (7.4,-1.9)
  {\textbf{Projected decoder}\\
   $(O,A,T_O,S_O)$};
\draw[flow] (sample.south -| uinput.north) -- (uinput.north);
\draw[flow] (sample.south -| oinput.north) -- (oinput.north);

\node[rule] (urule) at (0,-4.8)
  {Choose the fixed slice witness and\\its set of valid answers\\[2pt]
   $\mathcal A=(z+Z)\cap\{v:e_0^\top v=1\}$\\
   from this decoder's visible advice.\\[6pt]
   \textit{Private randomness:}\\
   sample $a_U$ uniformly from $\mathcal A$.};
\node[rule] (orule) at (7.4,-4.8)
  {Set $H=\ker A$. Choose $M_0$\\with $AM_0=S_O$.\\[5pt]
   \textit{Private randomness:}\\
   sample uniform $\sigma\in K^*$, then $\Phi$\\
   with probability $|\widehat f_\sigma(\Phi)|^2$.\\[5pt]
   If $\sigma|_H\ne0$, use a fixed $h\in H$\\
   with $\sigma(h)=1$; set $a_O=\Phi(h)$.};
\draw[flow] (uinput) -- (urule);
\draw[flow] (oinput) -- (orule);

\node[output] (uoutput) at (0,-7.9)
  {$a_U\in L_U,\qquad e_0^\top a_U=1$};
\node[output] (ooutput) at (7.4,-7.9)
  {$a_O\in L_O,\qquad e_0^\top a_O=1$};
\draw[flow] (urule) -- (uoutput);
\draw[flow] (orule) -- (ooutput);

\node[box,text width=8.2cm,minimum height=11mm] (verify) at (3.7,-9.6)
  {\textbf{Referee checks}\quad $\bar\pi a_U=a_O$};
\draw[flow] (uoutput.south) -- (verify.north -| uoutput.south);
\draw[flow] (ooutput.south) -- (verify.north -| ooutput.south);
\end{tikzpicture}
\caption{The two decoders use independent random tapes. Neither receives the
sampled $M_O$ or the projection $\bar\pi$; the right decoder chooses its own
representative $M_0$ from the visible advice. Every displayed advice map
includes its affine intercept, through the homogeneous coordinate. The
transferred slice is used only to analyze these private rules: the right
decoder does not use the left witness or answer set. If the left advice is
not good, or if $\sigma|_H=0$ on the right, that decoder returns an arbitrary
valid answer.}
\label{fig:advice}
\end{figure}

\FloatBarrier

\begin{lemma}[A transferred slice gives agreement]
\label{lem:fourier-coset-decoding}
Fix joint data satisfying Lemma~\ref{lem:projected-witness}.  Conditional
on these data, the two rules just defined, using independent randomness,
agree under $\bar\pi$ with probability at least
\begin{equation}
 2^{-\ell r_{\rm s}}(\alpha/8)^2\,2^{-r_{\rm s}}.
 \label{eq:conditional-decoding}
\end{equation}
\end{lemma}

\begin{proof}
Put $Z'=\bar\pi Z$, so $\dim Z'\le r_{\rm s}$.  The projected column
specifications in the row fiber fix the restriction $N|_{Z'}$ to a linear
map $b:Z'\to H$.  The resulting affine slice is nonempty; choose one
$N_*$ in it.  For every $\sigma\in K^*$ define
\[
 C_\sigma=\E\bigl[f_\sigma(N)(-1)^{\sigma(Nz')}
                       \mid N|_{Z'}=b\bigr].
\]
The character expansion of Equation~\eqref{eq:transferred-witness} gives
\[
 \E_\sigma\bigl[(-1)^{\sigma(M_0z'+u)}C_\sigma\bigr]
 \ge\alpha/4,
 \qquad\text{so}\qquad
 \E_\sigma|C_\sigma|\ge\alpha/4.
\]
Exactly a fraction $2^{-\dim H}$ of the characters have trivial
restriction to $H$.  Since $\dim H\ge\ell-r_{\rm s}$ and
$|C_\sigma|\le1$, Equation~\eqref{eq:matrix-parameters} implies
\begin{equation}
 \E_\sigma\bigl[\mathbf 1_{\{\sigma|_H\ne0\}}|C_\sigma|\bigr]
 \ge\alpha/8.
 \label{eq:nontrivial-correlation}
\end{equation}

The annihilator in $X^*$ of matrices vanishing on $Z'$ is $H^*\otimes Z'$.
For $\theta_\sigma=(\sigma|_H)\otimes z'$, Fourier expansion on the
affine slice therefore gives
\[
 C_\sigma
 =\sum_{\Phi\in\theta_\sigma+H^*\otimes Z'}
       \widehat f_\sigma(\Phi)
       (-1)^{\langle\Phi+\theta_\sigma,N_*\rangle}.
\]
This sum has $2^{\dim H\dim Z'}\le2^{\ell r_{\rm s}}$ terms.
Cauchy--Schwarz, first on each sum and then in the average over $\sigma$,
shows from Equation~\eqref{eq:nontrivial-correlation} that
\begin{align*}
 &\E_\sigma\left[
   \mathbf 1_{\{\sigma|_H\ne0\}}
   \sum_{\Phi\in\theta_\sigma+H^*\otimes Z'}
             |\widehat f_\sigma(\Phi)|^2\right]\\
 &\hspace{15mm}\ge
 2^{-\ell r_{\rm s}}
 \E_\sigma\bigl[\mathbf 1_{\{\sigma|_H\ne0\}}C_\sigma^2\bigr]
 \ge 2^{-\ell r_{\rm s}}(\alpha/8)^2.
\end{align*}
This is a lower bound on the probability that the $O$-decoder samples a
nontrivial restriction and a frequency in the indicated coset.

Every such frequency has the form
$\Phi=(\sigma|_H)\otimes z'+\Psi$ with $\Psi\in H^*\otimes Z'$.
For the decoder's choice of $h$ it satisfies
\[
 \Phi(h)=z'+\Psi(h)\in z'+Z'.
\]
As already proved, a frequency of positive mass gives first bit one.
Equation~\eqref{eq:projected-answer-set} therefore places the returned
vector in $\bar\pi\mathcal A$.  Independently, the $U$-decoder samples
uniformly from $\mathcal A$ and chooses a preimage with probability at
least $1/|\mathcal A|\ge2^{-r_{\rm s}}$.  Multiplying the bounds proves
Equation~\eqref{eq:conditional-decoding}.
\end{proof}

\begin{proof}[Proof of Proposition~\ref{prop:matrix-decoding}]
Use the two decoders above.  Lemma~\ref{lem:projected-witness} supplies a
set of joint data of probability at least $\gamma_1$, and
Lemma~\ref{lem:fourier-coset-decoding} gives the stated conditional
agreement on that set.  Thus one may take
\[
 \gamma=\gamma_1\,2^{-\ell r_{\rm s}}(\alpha/8)^2\,2^{-r_{\rm s}}>0.
\]
The largeness requirements for $k$ come only from
Theorem~\ref{thm:shortcode}, the concentration bound in
Lemma~\ref{lem:good-row-advice}, and the total variation estimate in
Lemma~\ref{lem:projected-tv}.  They are uniform in the parity instance
and table family.  Witness choices and Fourier sampling need not be efficient:
they define strategies in a finite game for the soundness argument.
Finally, averaging over the provers' independent random tapes also gives
deterministic strategies with at least the same agreement probability.
\end{proof}
